% Declaração de uso de ferramentas de IA generativa na elaboração do trabalho (tarefa 88).
% Regra do orientador (27/10): todo uso de IA na pesquisa é explicado e apontado no artigo, por etapa,
% na voz da equipe e sem coautoria. Nem o template da FCI nem as páginas públicas da Mackenzie pedem
% formato próprio; por isso, uma seção curta antes das referências.
\section*{Declaração de uso de inteligência artificial generativa}
\label{sec:declaracao-ia}

A equipe usou ferramentas de inteligência artificial (IA) generativa na elaboração deste trabalho, nas etapas descritas abaixo. Elas não são autoras: o que produziram foi revisado pela equipe, que responde pelo conteúdo. Os modelos de linguagem que fazem parte do sistema, os agentes da seção~\ref{sec:mcp}, são objeto de estudo, e não ferramentas de elaboração, e não entram nesta declaração.

Matheus Queiroz Gregorin usou o Claude Code, assistente de programação da Anthropic, em três etapas. No código, ele gerou os scripts das figuras e o comando que as regera, e o verificador do artigo contra as regras do template, com os testes automatizados. No texto, ele redigiu, a partir do esqueleto do artigo, trechos da Introdução, das seções 2.1 e 2.2, da Metodologia e dos Resultados, as Limitações, a lista de siglas e verbetes do glossário. Na análise, ele leu o artigo de \citeonline{neto2023} para a comparação com o baseline publicado e conferiu os números do texto contra os relatórios do repositório. Tudo passou por revisão em pull request de outro integrante, pelos testes automatizados e pela integração contínua, e os números foram conferidos contra os relatórios gerados pelos scripts do projeto \cite{oliveira2026}.

\pendente{ferramentas de IA generativa de Bernardo Souza Oliveira, José Victor Roling e Pedro Henrique Cagnoni Guimarães: quais, em que etapas e como cada um conferiu o que elas produziram (tarefa 88).}
